\section{Unmarked determinant scale}\label{c4:determinant}

Before estimating the marked radial sum we give the unmarked magnitude needed for comparison. Recall the determinant/color identities with \(P_0=2+\sqrt2,\ p_{\rm site}(t,s)=t^2+\sqrt2 ts+s^2,\ k(t,s)=ts/p_{\rm site}(t,s)\):
\[
 D_+(X)=P_0^N\prod_{i<j}[p_{\rm site}(t_i,t_j)/(t_i-t_j)]^2\,\det[k(t_i,t_j)],\qquad
 \tau_X=D_+(X)^2/\prod_{i<j}p_{\rm site}(t_i,t_j)(t_i^2+t_j^2).
\]
Writing \(y_i=\log t_i\), \(k=1/(2\cosh(y_i-y_j)+\sqrt2)\) has Fourier density 
\[\rho(T)=\sinh(\pi T/4)/(\sqrt2\sinh(\pi T))\]
 in \(\int e^{iT(y_i-y_j)}\rho(T)dT\). Indeed integrate \(e^{iyT}\sinh(a'T)/\sinh(\pi T)\) by residues for \(y>0,\ 0<a'<\pi\) in the upper half plane (upper edges at half-integers in imaginary part): this gives \(-2\sum_{n\ge1}(-1)^n\sin(a'n)e^{-ny}=\sin a'/(\cosh y+\cos a')\). Vandermonde division by confluence in rows and columns yields, at all \(t_i=1\),
\[
 D_+=P_0^{N^2}\mathcal D_N(\rho),\qquad
 \mathcal D_N(\rho)=\det\!\left[\int T^{j+k}\rho(T)dT/(j!k!)\right]_{j,k=0}^{N-1},
\]
in particular nonzero; the row/column \(i\)-powers from derivatives cancel. Set \(b'=3/4,\ \sigma(T)=\operatorname{sech}(\pi b'T)/(2\sqrt2)\) of mass \(u_0=1/(2\sqrt2 b')\), \(\sigma_0=\sigma/u_0,\ r(T)=\rho/\sigma=1-1/(2\cosh(\pi T/2))\). The corresponding Meixner--Pollaczek generating function~\cite[Equation~18.23.7]{DLMF} gives the following elementary orthogonality calculation. The orthonormal polynomials for \(\sigma_0\) are generated by
\[
 \sum q_n(T)z^n=(1+z^2)^{-1/2}\exp(2b'T\arctan z)
\]
because two such generating functions integrate as product to \(1/(1-zz')\); use \(\int e^{uT}\sigma_0(T)dT=1/\cos(u/(2b'))\). This integral follows by \(x=e^{2\pi b'T}\) using \(\int_0^\infty x^{s-1}/(1+x)\,dx=\pi/\sin(\pi s)\) for \(0<\Re s<1\) (integrate around the slit by residues). Thus \(\mathcal D_N(\sigma)=u_0^N(2b')^{-N(N-1)}\) by triangular leading coefficients.

\subsection{A reproducing-kernel comparison}

The explicit comparison weight gives the quadratic and linear exponential factors. We now calculate the remaining power of $N$ by interpolating between the comparison and physical moment determinants.

Let \(\mathcal K_{N,t}\) be the reproducing kernel of degree \(<N\) polynomials under \(\sigma_0 r^t\), \(0\le t\le1\). Then
\[
 \sigma_0(T)r(T)^t\mathcal K_{N,t}(T,T)/\log N\ \longrightarrow\ b'/\pi
\]
dominated by a polynomial in \(|T|\). For \(t=0\), put \(\alpha=1/2+i b'T\) just here. The generating function factors as \((1+i z)^{-\alpha}(1-i z)^{-\bar\alpha}\). Uniformly on compact real \(T\),
\[
 q_n(T)=n^{-1/2}[(-i)^n n^{i b'T}\mathcal A(T)+i^n n^{-i b'T}\overline{\mathcal A(T)}]+o(n^{-1/2}),
 \qquad \mathcal A(T)=2^{-\bar\alpha}/\Gamma(\alpha).
\]
Indeed subtract each leading singular branch with the other factor set to its value there; the remainder on the circle is continuous with uniformly integrable first derivative, giving coefficients \(O(1/n)\). Binomial coefficients have the displayed asymptotic (scale the beta integral \(B(\alpha,n)\)); \(|\Gamma(\alpha)|^2=\pi/\cosh(\pi b'T)\) by the same slit integral (beta identity by multiplying the two Euler gamma integrals and changing to sum and proportion coordinates). Consequently on bounded real \(v'\),
\[
 \mathcal K_{N,0}(T,T+v'/\log N)/\log N\longrightarrow
       2|\mathcal A(T)|^2\,\frac{\sin(b'v')}{b'v'} .
\]
The alternating product terms average to zero by partial summation, and harmonic averaging of continuous functions of \(\log n/\log N\) gives Lebesgue measure on \([0,1]\). In particular the probability law \(|\mathcal K_{N,0}(T,S)|^2\sigma_0(S)dS/\mathcal K_{N,0}(T,T)\) concentrates at \(T\): its mass on \(|S-T|\le R'/\log N\) tends to \(\int_{-R'}^{R'}(b'/\pi)(\sin(b'v')/(b'v'))^2dv'\), tending to 1 (e.g. Parseval). Testing the \(t=0\) kernel in the extremal evaluation norm for \(\mathcal K_{N,t}(T,T)\), and conversely reproducing an arbitrary test polynomial using \(\mathcal K_{N,0}\) and applying Cauchy--Schwarz, bounds \(\mathcal K_{N,t}/\mathcal K_{N,0}\) on diagonal between \(1/\mathbb E(r^t)\) and \(\mathbb E(r^{-t})\) with this law. This proves the limit since \(1/2\le r\le1\).

For domination, the circle generating function has magnitude \(\le C e^{\pi b' |T|/2}\) times angular distance to \(\{\pm i\}\) to power \(-1/2\), and derivative \(\le C(1+|T|) e^{\pi b'|T|/2}\) times this distance to power \(-3/2\) (\(|\Re\arctan z|\le\pi/4\) in the disk). Split its coefficient integral at distance \(1/(n+1)\) and integrate by parts on the complement. This gives \(|q_n(T)|\le C(1+|T|)e^{\pi b'|T|/2}(n+1)^{-1/2}\), sufficient since \(\mathcal K_{N,t}\le2\mathcal K_{N,0}\) on diagonal. Now \(\partial_t\log\mathcal D_N(\sigma r^t)=\int \mathcal K_{N,t}(T,T)\sigma_0(T)r(T)^t\log r(T)dT\) by the Gram derivative. Integrating in \(t\) with the exponential decay of \(\log r\) gives
\[
 (\log\mathcal D_N(\rho)-\log\mathcal D_N(\sigma))/\log N
       \longrightarrow (b'/\pi)\int\log r(T)dT=-7/24.
\]
For the integral use \(r=(1-e^{-s})(1+e^{-3s})/(1-e^{-4s}),\ s=\pi|T|/2\), expanding the logs (sum \(n^{-2}=\pi^2/6\) also by the sine product above). Thus
\[
 \tau_N=[P_0^2/((2b')^2\sqrt{2P_0})]^{N^2}[\sqrt{2P_0}/2]^N N^{-7/12+o(1)}.
\]
The two constants are respectively \(d_R(h)/d_C(h)\) and \(|w_-|\): by the given pairings 
\[d_C(h)=8^{-4/3}(\sqrt{2-\sqrt2})^{4/3}(\sqrt2)^{-1/3}(\sqrt{2+\sqrt2})^{-5/3}2^{-2/3},\]
\(d_R(h)=6^{-4/3}(\sqrt3)^{-4/3}2^{-4/3}\). Therefore the unmarked magnitude of \eqref{c4:Ang-4} divided by \(d_R(h)^{N^2}\) is \(N^{-7/12+o(1)}\).

